\section{Research Questions and Hypotheses}

\subsection{Research Questions}
This study is organized around one overarching question and four specific research questions that structure the experimental validation. The overarching question is:

\begin{quote}
\textbf{RQ0 (Overarching):} When Vietnamese legal knowledge is ingested \emph{continually}, which adaptation paradigm best avoids catastrophic forgetting across legal question answering (QA), natural language inference (NLI), and legal syllogism, and does a \emph{selective} knowledge store preserve retrieval sharpness where a \emph{flat} store degrades?
\end{quote}

This decomposes into four specific questions:
\begin{itemize}
    \item \textbf{Research Question 1 (RQ1 -- Forgetting under continual ingestion):} How severely does parametric continual learning (Continual Pre-Training / Continual Instruction Tuning) degrade on QA, NLI, and syllogism as the volume of ingested legal knowledge grows?
    \item \textbf{Research Question 2 (RQ2 -- RAG forgetting-resistance):} Does a frozen-base retrieval-augmented framework avoid weight-level forgetting and sustain (or improve) task performance as ingested knowledge grows, relative to a parametric baseline of the \emph{same} model size?
    \item \textbf{Research Question 3 (RQ3 -- Selective memory):} As the knowledge store grows, does selective memory (importance, decay, and compliance gating) preserve retrieval recall where a flat store loses it to accumulated index noise?
    \item \textbf{Research Question 4 (RQ4 -- Reasoning ceiling):} How does the benefit of retrieval vary with reasoning difficulty across the QA $\rightarrow$ NLI $\rightarrow$ syllogism progression?
\end{itemize}

\subsection{Research Hypotheses}
Corresponding to these questions, we formulate four empirically testable hypotheses:
\begin{itemize}
    \item \textbf{Hypothesis 1 (H1 -- Catastrophic forgetting under continual ingestion):} As the volume of continually ingested legal knowledge increases, parametric continual learning (CPT/CIT) degrades \emph{catastrophically}---a sharp drop in task performance and a strongly negative Backward Transfer---mirroring the collapse reported for continual post-training \parencite{su2026retrievablegradientscontinualposttraining}.
    \item \textbf{Hypothesis 2 (H2 -- RAG forgetting-resistance):} Under the same continual-ingestion schedule, the frozen-base RAG framework exhibits Backward Transfer close to zero and non-decreasing task performance, \emph{while a parametric baseline of identical model size degrades significantly}. The hypothesis is thus stated and tested as a \emph{measurable, falsifiable contrast} against that baseline---it can be refuted if the RAG framework also degrades (e.g.\ through retrieval-side interference as the store grows)---rather than as a trivial consequence of freezing the weights.
    \item \textbf{Hypothesis 3 (H3 -- Selective memory preserves recall):} As the knowledge store grows, a vanilla flat RAG store loses retrieval recall to accumulated noise, whereas selective memory---combining the legal importance score, Ebbinghaus-based decay, and the compliance hard-gate---maintains recall at the same active store size, keeping the index compact without breaching coverage of in-force law.
    \item \textbf{Hypothesis 4 (H4 -- Task-dependent reasoning ceiling):} The benefit of retrieval is \emph{task-dependent}: substantial on QA, moderate on NLI, and limited on syllogistic reasoning, where retrieval supplies the major premise but cannot substitute for the inference capacity of a small, frozen base model. Improvements therefore taper along the QA $\rightarrow$ NLI $\rightarrow$ syllogism progression.
\end{itemize}
